\chapter{Programs — Rules}
%%% generated by the blueprint pipeline from TCSlib.Programs.Light.Rules
\section{Overview}

Declarations of the Lean module \texttt{TCSlib.Programs.Light.Rules}.

\section{Declarations}

\begin{definition}[State update after a statement]
\lean{Light.Stmt.after}
\label{Light.Stmt.after}
\leanok
\uses{Light.Cond.Holds, Light.Expr.val, Light.State, Light.Stmt}
A state consists of an integer-valued location map and an integer-valued memory map,
both indexed by natural numbers. Define the state after executing a statement as
follows: a skip, procedure call, or while statement leaves the state unchanged; an
assignment updates the specified location with the expression’s value in the original
state, and a store updates the memory cell whose index is the natural-number conversion
of the address expression’s value, storing the value of the other expression. A sequence
applies its first statement and then its second, while a conditional applies the branch
selected by whether its condition holds in the original state.
\end{definition}

\begin{definition}[Statement block cost]
\lean{Light.Stmt.blockCost}
\label{Light.Stmt.blockCost}
\leanok
\uses{Light.Cond.cost, Light.Expr.cost, Light.Stmt}
A variable assignment has cost one plus the cost of its expression, and a memory store
has cost one plus the costs of its address and value expressions. Sequential composition
has the sum of the two statement costs; a conditional has cost one plus the condition
cost plus the larger branch cost. Every other statement has cost zero.
\end{definition}

\begin{definition}[Block safety]
\lean{Light.Stmt.BlockSafe}
\label{Light.Stmt.BlockSafe}
\leanok
\uses{Light.Cond.Holds, Light.Cond.Safe, Light.Expr.Safe, Light.Expr.val, Light.Limits, Light.Limits.Addr, Light.State, Light.Stmt, Light.Stmt.after}
Given limits and a state, a skip statement is safe. An assignment is safe when its
expression is safe; a store is safe when its address and value expressions are safe and
its address lies in the permitted range. A sequence is safe when its first statement is
safe in the initial state and its second is safe after the first statement; a
conditional is safe when its condition is safe and each branch is safe whenever that
branch is selected. While loops and procedure calls are not safe.
\end{definition}

\begin{definition}[Safe execution to a postcondition]
\lean{Light.Stmt.Runs}
\label{Light.Stmt.Runs}
\leanok
\uses{Light.Limits, Light.State, Light.Stmt, Light.Stmt.BlockSafe, Light.Stmt.after}
For a resource limit record $L$, a statement $s$, an initial state $\sigma$, and a
predicate $R$ on states, say that $s$ runs from $\sigma$ to $R$ under $L$ precisely when
$s$ is block-safe under $L$ in $\sigma$ and $R$ holds of the state obtained by executing
$s$ from $\sigma$.
\end{definition}

\begin{theorem}[Monotonicity of the postcondition]
\lean{Light.Stmt.Runs.mono}
\label{Light.Stmt.Runs.mono}
\leanok
\uses{Light.State, Light.Stmt, Light.Stmt.Runs}
\difficulty{1}
Let $L$ be resource limits, $s$ a statement, $\sigma$ a state, and $R,R'$ predicates on
states. If $s$ is block-safe from $\sigma$ and its resulting state satisfies $R$, and
every state satisfying $R$ also satisfies $R'$, then $s$ is block-safe from $\sigma$ and
its resulting state satisfies $R'$.
\end{theorem}

\begin{theorem}[Sequential composition preserves runs]
\lean{Light.Stmt.Runs.seq}
\label{Light.Stmt.Runs.seq}
\leanok
\uses{Light.State, Light.Stmt, Light.Stmt.Runs}
\difficulty{1}
Let $L$ be resource limits, let $s,t$ be statements, let $\sigma$ be a state, and let
$R$ be a predicate on states. If $s$ is block-safe from $\sigma$ and its resulting state
is one from which $t$ is block-safe and reaches a state satisfying $R$, then the
sequential composition of $s$ followed by $t$ is block-safe from $\sigma$ and its
resulting state satisfies $R$.
\end{theorem}

\begin{theorem}[True branch of a conditional run]
\lean{Light.Stmt.Runs.ite_pos}
\label{Light.Stmt.Runs.ite_pos}
\leanok
\uses{Light.Cond, Light.Cond.Holds, Light.Cond.Safe, Light.State, Light.Stmt, Light.Stmt.Runs, Light.Stmt.after}
\difficulty{3}
Let $L$ be resource limits, $c$ a condition, $s,t$ statements, $\sigma$ a state, and $R$
a predicate on states. If $s$ is block-safe from $\sigma$ and its resulting state
satisfies $R$, if $c$ holds in $\sigma$, and if $c$ is safe under $L$ in $\sigma$, then
the conditional with branches $s,t$ is block-safe from $\sigma$ and its resulting state
satisfies $R$.
\end{theorem}

\begin{theorem}[False branch of a conditional run]
\lean{Light.Stmt.Runs.ite_neg}
\label{Light.Stmt.Runs.ite_neg}
\leanok
\uses{Light.Cond, Light.Cond.Holds, Light.Cond.Safe, Light.State, Light.Stmt, Light.Stmt.Runs, Light.Stmt.after}
\difficulty{3}
Let $L$ be resource limits, $c$ a condition, $s,t$ statements, $\sigma$ a state, and $R$
a predicate on states. If $t$ is block-safe from $\sigma$ and its resulting state
satisfies $R$, if $c$ does not hold in $\sigma$, and if $c$ is safe under $L$ in
$\sigma$, then the conditional with branches $s,t$ is block-safe from $\sigma$ and its
resulting state satisfies $R$.
\end{theorem}

\begin{theorem}[Execution of a safe straight-line block]
\lean{Light.Ends.of_blockSafe}
\label{Light.Ends.of_blockSafe}
\leanok
\uses{Light.Ends, Light.Ends.ite, Light.Ends.seq, Light.Ends.set, Light.Ends.skip, Light.Ends.store, Light.State, Light.Stmt, Light.Stmt.BlockSafe, Light.Stmt.after, Light.Stmt.blockCost}
\difficulty{5}
For any statement $s$, initial state $\sigma$, cost bound $T$, and state predicate $Q$,
if $s$ is block-safe from $\sigma$, its block cost is at most $T$, and $Q$ holds of the
state obtained by applying $s$'s prescribed state transformation to $\sigma$, then $s$
has an execution from $\sigma$ with cost at most $T$ ending in a state satisfying $Q$.
\end{theorem}

\begin{theorem}[Execution from a safe block]
\lean{Light.Ends.block}
\label{Light.Ends.block}
\leanok
\uses{Light.Ends, Light.Ends.of_blockSafe, Light.State, Light.Stmt, Light.Stmt.Runs, Light.Stmt.blockCost}
\difficulty{1}
Let $s$ be a statement and $\sigma$ a state. Suppose every operation in the block is
safe in sequence from $\sigma$, and the state obtained by applying the block's
assignments and stores satisfies $Q$. If the block cost is at most $T$, then $s$ has an
execution from $\sigma$ of cost at most $T$ ending in a state satisfying $Q$.
\end{theorem}

\begin{theorem}[Sequential composition of bounded executions]
\lean{Light.Ends.next}
\label{Light.Ends.next}
\leanok
\uses{Light.Ends, Light.Ends.seq, Light.State, Light.Stmt}
\difficulty{2}
Let $T,T_1$ be natural numbers, let $s_1,s_2$ be statements, and let $Q$ be a predicate
on states. If $s_1$ has an execution from $\sigma$ with cost at most $T_1$ such that its
final state admits an execution of $s_2$ with cost at most $T-T_1$ ending in a state
satisfying $Q$, and $T_1 \le T$, then the sequential composition of $s_1$ followed by
$s_2$ has an execution from $\sigma$ with cost at most $T$ ending in a state satisfying
$Q$.
\end{theorem}

\begin{theorem}[Associativity of bounded sequential execution]
\lean{Light.Ends.seqAssoc}
\label{Light.Ends.seqAssoc}
\leanok
\uses{Light.Ends, Light.State, Light.Stmt}
\difficulty{4}
Let $lim$ be resource limits, $P$ a program, $d$ a call depth, and $s_1,s_2,s_3$
statements. For any state $sigma$, time bound $T$, and postcondition $Q$, if
$s_1;(s_2;s_3)$ ends from $sigma$ within time $T$ satisfying $Q$, then $(s_1;s_2);s_3$
also ends from $sigma$ within time $T$ satisfying $Q$.
\end{theorem}

\begin{theorem}[Skip before a program]
\lean{Light.Ends.skipThen}
\label{Light.Ends.skipThen}
\leanok
\uses{Light.Ends, Light.Ends.seq, Light.Ends.skip, Light.State, Light.Stmt}
\difficulty{2}
For any limit record, program, execution depth, state $\sigma$, budget $T$, statement
$s$, and postcondition $Q$, if $s$ can execute from $\sigma$ with cost at most $T$ and
end in a state satisfying $Q$, then the sequence consisting of skip followed by $s$ can
also execute from $\sigma$ with cost at most $T$ and end in a state satisfying $Q$.
\end{theorem}

\begin{theorem}[Removing a final skip]
\lean{Light.Ends.skipLast}
\label{Light.Ends.skipLast}
\leanok
\uses{Light.Ends, Light.State, Light.Stmt}
\difficulty{4}
For any limit record, program, execution depth, state $\sigma$, budget $T$, statement
$s$, and postcondition $Q$, if the sequence consisting of $s$ followed by skip can
execute from $\sigma$ with cost at most $T$ and end in a state satisfying $Q$, then $s$
itself can execute from $\sigma$ with cost at most $T$ and end in a state satisfying
$Q$.
\end{theorem}

\begin{theorem}[Reflexivity of bounded sequential execution]
\lean{Light.Ends.seqSelf}
\label{Light.Ends.seqSelf}
\leanok
\uses{Light.Ends, Light.State, Light.Stmt}
\difficulty{1}
Let $lim$ be resource limits, $P$ a program, $d$ a call depth, and $s_1,s_2$
statements. For any state $sigma$, time bound $T$, and postcondition $Q$, if $s_1;s_2$
ends from $sigma$ within time $T$ satisfying $Q$, then $s_1;s_2$ ends from $sigma$
within time $T$ satisfying $Q$.
\end{theorem}

\begin{theorem}[Sequential composition from an intermediate postcondition]
\lean{Light.Ends.pieceThen}
\label{Light.Ends.pieceThen}
\leanok
\uses{Light.Ends, Light.Ends.mono, Light.Ends.next, Light.State, Light.Stmt}
\difficulty{3}
Let $s_1,s_2$ be statements, $T,T_1$ natural numbers, and $R,Q$ state predicates. If
$s_1$ has an execution from $\sigma$ of cost at most $T_1$ ending in a state satisfying
$R$, and from every such state $s_2$ has an execution of cost at most $T-T_1$ ending in
a state satisfying $Q$, with $T_1\le T$, then the sequential composition of $s_1$
followed by $s_2$ has an execution from $\sigma$ with cost at most $T$ ending in a state
satisfying $Q$.
\end{theorem}

\begin{theorem}[Weakening an execution postcondition]
\lean{Light.Ends.pieceLast}
\label{Light.Ends.pieceLast}
\leanok
\uses{Light.Ends, Light.Ends.mono, Light.State, Light.Stmt}
\difficulty{1}
Let $s$ be a statement, $\sigma$ an initial state, $T_1\le T$ natural-number cost
bounds, and $R,Q$ state predicates. If $s$ has an execution from $\sigma$ with cost at
most $T_1$ ending in a state satisfying $R$, and every state satisfying $R$ also
satisfies $Q$, then $s$ has an execution from $\sigma$ with cost at most $T$ ending in a
state satisfying $Q$.
\end{theorem}

\begin{theorem}[Bounded execution of an assignment]
\lean{Light.Ends.setLast}
\label{Light.Ends.setLast}
\leanok
\uses{Light.Ends, Light.Ends.set, Light.Expr, Light.Expr.Safe, Light.Expr.cost, Light.Expr.val, Light.State}
\difficulty{1}
Let $\mathrm{lim}$ be resource limits, $P$ a program, and $d$ a call depth. For
local-value and memory maps $\mathrm{loc},\mu:\mathbb{N}\to\mathbb{Z}$, an index $x$, an
expression $e$, a time bound $T$, and a postcondition $Q$, suppose $Q$ holds of the
state obtained by updating $\mathrm{loc}$ at $x$ to the value of $e$ in the state
$(\mathrm{loc},\mu)$ while leaving memory equal to $\mu$. If $e$ is safe in
$(\mathrm{loc},\mu)$ and its cost plus one is at most $T$, then assigning $e$ to $x$
from $(\mathrm{loc},\mu)$ ends within time $T$ in a state satisfying $Q$.
\end{theorem}

\begin{theorem}[Assignment followed by bounded execution]
\lean{Light.Ends.setThen}
\label{Light.Ends.setThen}
\leanok
\uses{Light.Ends, Light.Ends.next, Light.Ends.set, Light.Expr, Light.Expr.Safe, Light.Expr.cost, Light.Expr.val, Light.State, Light.Stmt}
\difficulty{1}
Let $L$ be resource limits, $P$ a program, $d$ a call depth, and
$	ext{loc},	ext{mem}:bn	obz$ an initial state. For an index $x$, an expression $e$, a
statement $s$, a time bound $T$, and a postcondition $Q$, suppose that executing $s$
from the state obtained by assigning the value of $e$ in the initial state to local $x$
ends within $	ext{max}(T-(	ext{cost}(e)+1),0)$ steps and satisfies $Q$. If $e$ is safe
in the initial state and $	ext{cost}(e)+1\le T$, then executing the assignment of $e$ to
$x$ followed by $s$ from the initial state ends within $T$ steps and satisfies $Q$.
\end{theorem}

\begin{theorem}[Store statement with final state]
\lean{Light.Ends.storeLast}
\label{Light.Ends.storeLast}
\leanok
\uses{Light.Ends, Light.Ends.store, Light.Expr, Light.Expr.Safe, Light.Expr.cost, Light.Expr.val, Light.Limits.Addr, Light.State}
\difficulty{1}
Let a limit record, program, execution depth, local function $\mathrm{loc}$, memory
function $\mu$, budget $T$, address expression $a$, value expression $e$, and
postcondition $Q$ be given. Suppose $Q$ holds of the state with local function
$\mathrm{loc}$ and memory updated at the natural-number conversion of $a$'s value in
$(\mathrm{loc},\mu)$ to $e$'s value in that state. Suppose also that $a$ and $e$ are
safe in $(\mathrm{loc},\mu)$ and that the value of $a$ is a valid address, meaning it is
nonnegative and less than the space bound. If $\mathrm{cost}(a)+\mathrm{cost}(e)+1\le
T$, then the store statement can execute from $(\mathrm{loc},\mu)$ with cost at most $T$
and end in a state satisfying $Q$.
\end{theorem}

\begin{theorem}[Store followed by a bounded computation]
\lean{Light.Ends.storeThen}
\label{Light.Ends.storeThen}
\leanok
\uses{Light.Ends, Light.Ends.next, Light.Ends.storeLast, Light.Expr, Light.Expr.Safe, Light.Expr.cost, Light.Expr.val, Light.Limits.Addr, Light.State, Light.Stmt}
\difficulty{1}
Let a program, depth bound, initial location map $\ell:\mathbb N\to\mathbb Z$, memory
map $\mu:\mathbb N\to\mathbb Z$, expressions $a,e$, statement $s$, time bound
$T\in\mathbb N$, and postcondition $Q$ be given. If $s$ can execute from the state with
locations $\ell$ and memory updated at address $a(\ell,\mu)$ to value $e(\ell,\mu)$,
within time $T-(\operatorname{cost}(a)+\operatorname{cost}(e)+1)$, and end in a state
satisfying $Q$, and if $a,e$ are safe in the original state and $a(\ell,\mu)$ is a valid
address, then storing $e$ at address $a$ and executing $s$ from the original state ends
in a state satisfying $Q$ within time $T$, provided
$\operatorname{cost}(a)+\operatorname{cost}(e)+1\le T$.
\end{theorem}

\begin{theorem}[Bounded execution of a conditional]
\lean{Light.Ends.iteLast}
\label{Light.Ends.iteLast}
\leanok
\uses{Light.Cond, Light.Cond.Holds, Light.Cond.Safe, Light.Cond.cost, Light.Ends, Light.Ends.ite, Light.State, Light.Stmt}
\difficulty{2}
Fix a program environment, resource limits, and execution depth. Let $c$ be a condition
safe in an initial state $\sigma$, and let $s_1,s_2$ be statements. Suppose that if $c$
holds in $\sigma$, then $s_1$ can execute from $\sigma$ with cost at most
$T-(\mathrm{cost}(c)+1)$ to a state satisfying $Q$; suppose likewise that if $c$ does
not hold, then $s_2$ can execute with that cost bound to a state satisfying $Q$. If
$\mathrm{cost}(c)+1\le T$, the conditional with branches $s_1,s_2$ can execute from
$\sigma$ with cost at most $T$ to a state satisfying $Q$.
\end{theorem}

\begin{theorem}[Conditional followed by a continuation]
\lean{Light.Ends.iteThen}
\label{Light.Ends.iteThen}
\leanok
\uses{Light.Cond, Light.Cond.Holds, Light.Cond.Safe, Light.Cond.cost, Light.Ends, Light.State, Light.Stmt}
\difficulty{4}
Fix a program environment, resource limits, and execution depth. Let $c$ be a condition
and let $s_1,s_2,s$ be statements. In an initial state $\sigma$, suppose that when $c$
holds, the sequence of $s_1$ followed by $s$ can execute with cost at most
$T-(\mathrm{cost}(c)+1)$ to a state satisfying $Q$, and that when $c$ does not hold, the
corresponding sequence using $s_2$ can do so. If $c$ is safe in $\sigma$ and
$\mathrm{cost}(c)+1\le T$, then the conditional with branches $s_1,s_2$, followed by
$s$, can execute from $\sigma$ with cost at most $T$ to a state satisfying $Q$.
\end{theorem}

\begin{theorem}[Conditional followed by a continuation]
\lean{Light.Ends.iteIffThen}
\label{Light.Ends.iteIffThen}
\leanok
\uses{Light.Cond, Light.Cond.Holds, Light.Cond.Safe, Light.Cond.cost, Light.Ends, Light.Ends.iteThen, Light.State, Light.Stmt}
\difficulty{1}
Fix a program environment, resource limits, and execution depth. Let $c$ be a condition,
let $s_1,s_2,s$ be statements, and let $p$ be a proposition. In an initial state
$\sigma$, suppose that if $p$ holds then executing $s_1$ followed by $s$ can finish with
cost at most $T-(\mathrm{cost}(c)+1)$ in a state satisfying $Q$, and if $p$ does not
hold then the same is true for $s_2$ followed by $s$. Suppose $c$ is safe in $\sigma$
and holds there exactly when $p$ holds, and that $\mathrm{cost}(c)+1\le T$. Then
executing the conditional with branches $s_1,s_2$, followed by $s$, from $\sigma$ can
finish with cost at most $T$ in a state satisfying $Q$.
\end{theorem}

\begin{theorem}[Bounded execution of a while loop]
\lean{Light.Ends.whileBlock}
\label{Light.Ends.whileBlock}
\leanok
\uses{Light.Cond, Light.Cond.Holds, Light.Cond.Safe, Light.Cond.cost, Light.Ends, Light.Ends.block, Light.Ends.whileConst, Light.State, Light.Stmt, Light.Stmt.Runs, Light.Stmt.blockCost}
\difficulty{1}
Let $L$ be a resource limit with word bound $w\in\mathbb Z$, space bound $s\in\mathbb
N$, and depth bound $h\in\mathbb N$; let $P$ be a finite program, $d,n,T\in\mathbb N$,
and let $\sigma$ be an initial state, $c$ a loop condition, $b$ a body statement, $Q$ a
predicate on states, and $I_j$ a state predicate for each $j\in\mathbb N$. Suppose
$I_0(\sigma)$ holds, and for every $j<n$ and state $\tau$ with $I_j(\tau)$, the
condition $c$ is safe at $\tau$ and true there, while $b$ is block-safe at $\tau$ and
its resulting state satisfies $I_{j+1}$. Suppose also that for every state $\tau$
satisfying $I_n(\tau)$, $c$ is safe at $\tau$ and false there, and $Q(\tau)$ holds. If
$n\bigl(\operatorname{cost}(c)+1+\operatorname{blockCost}(b)\bigr)+\operatorname{cost}(c)+1\le
T$, then executing the loop with condition $c$ and body $b$ from $\sigma$, under limits
$L$, program $P$, and depth parameter $d$, ends in some state satisfying $Q$ with
execution cost at most $T$.
\end{theorem}

\begin{definition}[For-loop statement]
\lean{Light.Stmt.for}
\label{Light.Stmt.for}
\leanok
\uses{Light.Expr, Light.Stmt, Light.k, Light.v}
For a natural-number variable $i$, an upper-bound expression $h$, and a body statement,
define the for-loop statement to initialize $i$ to zero and then repeatedly execute the
body followed by incrementing $i$ by one, while $i<h$.
\end{definition}

\begin{theorem}[Bounded execution of a for-loop]
\lean{Light.Ends.for}
\label{Light.Ends.for}
\leanok
\uses{Light.Cond.cost, Light.Ends, Light.Ends.mono, Light.Ends.seq, Light.Ends.set, Light.Ends.whileConst, Light.Expr, Light.Expr.Safe, Light.Expr.cost, Light.Expr.val, Light.State, Light.Stmt, Light.Stmt.for, Light.k, Light.v}
\difficulty{6}
Let $	ext{lim}$ be execution limits, let $P$ be a program, and let $i,n,T$ be natural
numbers, with initial state $sigma$, upper-bound expression $hi$, loop body $body$, and
postcondition $Q$. Let $I(j,sigma)$ be an invariant. Suppose after setting local
variable $i$ to $0$, $I(0)$ holds. For every $j<n$, if a state has local value $i=j$ and
satisfies $I(j)$, assume the body executes within cost $b$, preserves the loop index at
$j$, and leaves a state satisfying $I(j+1)$ after that index is incremented by one.
Assume that at index $n$, $I(n)$ and local value $i=n$ imply $Q$; and whenever $j≤n$,
local value $i=j$ and $I(j)$ hold, the upper-bound expression is safe and evaluates to
$n$. If $n$ is at most the word limit and $n(	ext{cost}(hi)+b+7)+	ext{cost}(hi)+5≤T$,
then the for-loop terminates from $sigma$ within cost $T$ in a state satisfying $Q$.
\end{theorem}

\begin{theorem}[For-loop with an invariant on memory]
\lean{Light.Ends.forMem}
\label{Light.Ends.forMem}
\leanok
\uses{Light.Ends, Light.Ends.for, Light.Ends.mono, Light.Expr, Light.Expr.Safe, Light.Expr.cost, Light.Expr.val, Light.State, Light.Stmt, Light.Stmt.for}
\difficulty{4}
Let $	ext{lim}$ be execution limits, let $P$ be a program, and let $i,n,T,b$ be natural
numbers, with fixed local-variable function $loc$, initial memory $mu$, upper-bound
expression $hi$, loop body $body$, and postcondition $Q$. Let $I(j,mu')$ be a memory
invariant. Suppose $I(0,mu)$ holds. For each $j<n$, assume that from local variables
obtained by setting $i$ to $j$ in $loc$ and any memory satisfying $I(j)$, the body
executes within cost $b$, preserves the local-variable function, and leaves memory
satisfying $I(j+1)$. Assume that $I(n,mu')$ implies $Q$ for the state with local
variable $i$ set to $n$ and memory $mu'$. Also assume that whenever $j≤n$ and
$I(j,mu')$ holds, $hi$ is safe and evaluates to $n$ in that state. If $n$ is at most
the word limit and $n(	ext{cost}(hi)+b+7)+	ext{cost}(hi)+5≤T$, then the for-loop
terminates from $(loc,mu)$ within cost $T$ in a state satisfying $Q$.
\end{theorem}
